\documentclass[9pt,twocolumn,twoside]{pnas-new}
% Use the lineno option to display guide line numbers if required.
\setlength{\affilsep}{2pt}
\articletype{BIOLOGICAL SCIENCES: Neuroscience}

\templatetype{pnasresearcharticle} % Choose template
% {pnasresearcharticle} = Template for a two-column research article
% {pnasmathematics} %= Template for a one-column mathematics article
% {pnasinvited} %= Template for a PNAS invited submission


\begin{document}


\title{Second-order directed embedding and functional connectivity variability in larval zebrafish, with comparative analyses in two invertebrate systems}

% Use letters for affiliations, numbers to show equal authorship (if applicable) and to indicate the corresponding author
\author[a,b]{Dongmyeong Lee}
\author[a,c]{Euisun Kim} 
\author[d]{Richard Rosch} 
\author[a,b,c,e,1]{Hae-Jeong Park}

\affil[a]{Center for Systems Brain Science, Institute of Human Complexity and Systems Science, Yonsei University, Seoul 06130, Republic of Korea}
\affil[b]{Department of Nuclear Medicine, Department of Psychiatry, Yonsei University College of Medicine, Seoul 06130, Republic of Korea}
\affil[c]{BK21 PLUS Project for Medical Science, Yonsei University College of Medicine, Seoul 06130, Republic of Korea}
\affil[d]{Basic and Clinical Neuroscience Department, Institute of Psychiatry, Psychology and Neuroscience, King's College London}
\affil[e]{Epilepsy Research Institute, Brain Research Institute, Yonsei University College of Medicine, Seoul 06130, Republic of Korea}
% Please give the surname of the lead author for the running footer
\leadauthor{Lee}

% Please add a significance statement to explain the relevance of your work
\significancestatement{
Brain circuits must change their interactions flexibly from moment to moment. Most connectome studies emphasize strong connections, but this does not explain why circuits continually shift partners. Using larval zebrafish brain-wide calcium imaging and synapse-informed directed connectivity, we identify a cascade-like structural signature in which output-biased neurons target downstream neurons positioned to propagate signals. Across regions, stronger cascade embedding is linked to higher functional connectivity variability, distinguishing high-variability telencephalic regimes from other divisions; outgoing-projection perturbations reduce variability more than incoming ones in simulations, and a minimal directed model shows directional asymmetry can amplify variability. 
Comparative analyses in \textit{C.~elegans} and \textit{Drosophila} revealed related, but scale- and circuit-dependent, associations between high-FCV groups and directed embedding profiles.
}


% Please include corresponding author, author contribution and author declaration information
\authorcontributions{
D.L. and H.J.P. designed research; D.L. and E.K. performed research; R.R. contributed the calcium analysis; and D.L. and H.J.P. wrote the paper.}
\authordeclaration{The authors declare no competing interest.}
%\equalauthors{\textsuperscript{1}A.O.(Author One) contributed equally to this work with A.T. (Author Two) (remove if not applicable).}
\correspondingauthor{\textsuperscript{1}To whom correspondence should be addressed. E-mail: hjpark0@gmail.com}

% At least three keywords are required at submission. Please provide three to five keywords, separated by the pipe symbol.
\keywords{cascaded directional embedding | functional connectivity variability | directional coreness asymmetry | cross-species connectome comparison | network asymmetry}

\begin{abstract}
Some neural circuits repeatedly engage the same functional partners, whereas others continually reconfigure their interactions over time.
The synaptic architecture supporting this flexibility remains unclear.
Combining whole-brain calcium imaging with a synapse-informed directed structural network in larval zebrafish, we asked whether flexible functional coupling is related to second-order directed embedding: not only whether a circuit sends outputs, but whether those outputs reach downstream populations positioned to propagate activity further.
Regions with more variable functional interactions showed stronger output-side cascade embedding, forming a cascade-like wiring pattern.
This association was strongest for telencephalic circuits, which showed high temporal reconfiguration together with downstream output-biased embedding.
In simulations, disrupting outgoing projections reduced functional variability more strongly than disrupting incoming projections, and a minimal directed network showed that directional asymmetry can amplify fluctuations in functional coupling.
Comparative analyses in \textit{C.~elegans} and adult \textit{Drosophila} identified related group-level or region-level patterns at the finest scales where structural and functional data could be matched, although the dominant output-side or input-side expression differed across species.
Together, these results support cascaded directional embedding as a candidate structural correlate of temporally flexible circuit dynamics.
\end{abstract}



\dates{This manuscript was compiled on \today}
% \doi{\url{www.pnas.org/cgi/doi/10.1073/pnas.XXXXXXXXXX}}


\maketitle
\thispagestyle{firststyle}
\ifthenelse{\boolean{shortarticle}}{\ifthenelse{\boolean{singlecolumn}}{\abscontentformatted}{\abscontent}}{}

\Firstpage

%The \Firstpage command is used to format the first page text column size. The same size will be maintained for subsequent paragraph until the \Endparasplit or \Parasplit command is encountered.

Neural circuits differ not only in how strongly their components are coupled, but in how often those couplings change. Some circuits maintain stable functional relationships, whereas others continually reconfigure which partners they interact with from moment to moment---a distinction especially relevant to telencephalic/pallial circuits that support flexible cognition, goal-directed behavior, spatial representation, contextual adaptation, and associative plasticity \citep{miller2001integrative,duncan2001adaptive,sakai2006prefrontal,cole2013multi,braun2015dynamic,Yang2024SpatialCode,Dempsey2022ZebrafishMemorySynapses}. This temporal reconfiguration is captured by functional connectivity variability (FCV): the degree to which pairwise functional coupling fluctuates across time, reflecting a dynamical regime in which coordinated activity patterns shift transiently between configurations \citep{werner2007metastability,deco2017hierarchy,tognoli2014metastable}. Although structural connectomes can predict the average strength of functional coupling, they do not explain why some circuits are dynamically flexible while others remain stable, and this question has not been systematically tested in a whole-brain setting with cell-resolved directed structural information. What wiring motif is associated with functional interactions that reconfigure over time?

Structural connectivity (SC) constrains the strength and stability of functional connectivity (FC), as demonstrated across species and imaging modalities \citep{honey2009predicting,honey2010can,goni2014resting,deco2011emerging,rosch2024spontaneous}. These approaches, however, have focused almost exclusively on predicting the time-averaged strength of FC, relying on undirected or coarse-grained connectomes and emphasizing classical architectural features---modularity, clustering, and global efficiency---that quantify network segregation and integration. Whether and how structural connectivity accounts for the \emph{temporal variability} of functional interactions, rather than their mean strength, has received far less attention. \Endparasplit

More recent work highlights the role of directionality: asymmetric connectivity---imbalances between incoming and outgoing connections---can transiently amplify activity and facilitate transitions between network states \citep{battaglia2012dynamic,dahmen2025heterogeneity,murphy2009balanced,goldman2009memory,hennequin2014optimal}. These studies characterize directionality as a property of individual nodes, asking whether a given neuron broadcasts more than it integrates---a first-order description of the network. What has not been addressed is a second-order question: when an output-biased circuit sends projections outward, what is the directed role of the neurons it reaches? The same first-order output bias can be embedded in different downstream architectures depending on whether its targets are themselves output-biased or input-biased. Understanding this second-order structure is essential for linking directed synaptic architecture to the temporal dynamics of brain activity.


The dynamical effects of directionality should therefore depend not only on whether individual neurons are output- or input-dominated, but also on how these neurons are connected to one another \citep{ratnarajah2013structural,seguin2019inferring,harriger2012rich,esposito2014measuring,hadjiabadi2021maximally}. If broadcast-capable neurons preferentially target other broadcast-capable neurons, output bias is embedded in an outgoing cascade-like architecture. More generally, flexible dynamics may depend on the directed role of the populations reached by, or projecting into, a given circuit. We refer to this second-order arrangement as \emph{cascaded directional embedding}: a structural motif in which functional variability depends on the directed role of connected upstream or downstream populations.

This idea differs from the simpler view that circuits become more variable whenever they contain more output-biased elements. Instead, it predicts that variability depends on the broader directed context in which output- and input-biased elements are embedded. In some circuits, flexible dynamics may be shaped by where a region sends its outputs; in others, it may be shaped by the structure of the upstream inputs it receives. 

This framing leads to two hypotheses. First, FC variability should be associated with second-order directional embedding rather than simple output bias alone, but the relevant side of the cascade may differ across circuit architectures. Second, cascade topology should help distinguish broad functional regimes, such as telencephalic versus non-telencephalic or associative versus sensorimotor circuits.


Testing these hypotheses requires two capabilities that are rarely available together: synapse-informed cell-level structural information and whole-brain functional dynamics at cellular resolution. Defining cascade embedding within a regional subnetwork requires directionality resolved at the cellular projection level: the fiber-tract connectomes used in humans and non-human primates cannot identify which neurons within a region are output- or input-dominated and therefore cannot define directional coreness asymmetry at the subnetwork level. Testing its relationship to FC variability additionally requires a whole-brain functional phenotype spanning all major circuit divisions simultaneously, so that telencephalic and non-telencephalic circuits can be compared on equal footing. Larval zebrafish are well suited to this combination of requirements \citep{kinkhabwala2011structural,ahrens2013whole,cong2017rapid,kim2017pan,chen2018brain,kunst2019cellular}: whole-brain calcium imaging provides near-complete coverage of neuronal activity across all major brain divisions \citep{chen2018brain}, while cellular-resolution anatomical reconstructions supply synapse-informed directed structural connectivity \citep{kunst2019cellular,hildebrand2017wholebrain}.

Here, we test these hypotheses by combining whole-brain calcium imaging with a synapse-informed directed structural network in larval zebrafish, supplemented by large-scale biophysical simulations and topology-preserving rewiring perturbations.
We show that zebrafish FCV is associated with output-side second-order cascade topology rather than simple output bias alone.
In direction-specific rewiring, outgoing perturbations reduced FCV more strongly than incoming perturbations, and a minimal directed network model showed that directional asymmetry can amplify functional reconfiguration.
Comparative analyses in \textit{C.~elegans} and adult \textit{Drosophila} identified related group-level or region-level patterns at the finest available matched scales, suggesting that directed embedding may be relevant beyond zebrafish while remaining species- and scale-dependent. Together, these results support cascaded directional embedding as a candidate structural correlate of temporally flexible circuit dynamics.


\section*{Results}

\begin{figure*}[!ht]
\includegraphics[width=\linewidth]{figure/figure9_final.png}
\caption{
\textbf{Region-wise FC-dynamics and directed-interaction profiles reveal telencephalic output-drive organization.}
(A) Hierarchical clustering of zebrafish regional FC-dynamics and directed-interaction measures.
Rows show FC variability (FCV), functional connectivity strength (FCS), FC partner reconfiguration, net transfer entropy ($\mathrm{TE}_{\mathrm{net}}$), and neighborhood $\mathrm{TE}_{\mathrm{net}}$.
Values are the zebrafish five-measure regional summaries used in Fig.~1, after within-recording normalization and node-level averaging.
Columns indicate anatomical regions ordered by Ward hierarchical clustering.
The top annotation bar indicates four anatomical divisions: telencephalon (Tel), diencephalon (Di), mesencephalon (Mes), and hindbrain (Hind).
(B) Regional directed interaction network reconstructed from zebrafish regional-level FC and transfer-entropy data.
Displayed edges were restricted to regional pairs with FC-neighbor support in at least 4 of 7 subjects, using the stored BH-FDR-corrected regional-level FC-neighbor mask ($q < 0.001$).
Among FC-supported directed pairs with positive regional $\mathrm{TE}_{\mathrm{net}}$, only the top quartile was displayed.
Edge color indicates the source anatomical division. Node size and color intensity indicate rank-normalized FCV.
(C--G) Division-level distributions of the same five measures shown in panel A: FCV (C), FCS (D), FC partner reconfiguration (E), $\mathrm{TE}_{\mathrm{net}}$ (F), and neighborhood $\mathrm{TE}_{\mathrm{net}}$ (G).
Points indicate recording-by-region observations after within-recording z-scoring; box plots show median and interquartile range.
Regions are grouped into Tel, Di, Mes, and Hind. Group differences were assessed using Kruskal--Wallis tests followed by Holm-corrected two-sided Mann--Whitney U tests.
Kruskal--Wallis results were: FCV, $H = 76.847$, $p = 1.46\times 10^{-16}$, $\eta^2_H = 0.166$; FCS, $H = 109.288$, $p = 1.56\times 10^{-23}$, $\eta^2_H = 0.238$; FC partner reconfiguration, $H = 90.553$, $p = 1.67\times 10^{-19}$, $\eta^2_H = 0.196$; $\mathrm{TE}_{\mathrm{net}}$, $H = 39.246$, $p = 1.54\times 10^{-8}$, $\eta^2_H = 0.121$; and neighborhood $\mathrm{TE}_{\mathrm{net}}$, $H = 26.755$, $p = 6.62\times 10^{-6}$, $\eta^2_H = 0.079$.
Statistical significance is indicated as *$p < 0.05$, **$p < 0.01$, and ***$p < 0.001$.
}
\label{fig:fc_analysis}
\end{figure*}



\subsection*{Telencephalic circuits exhibit elevated FC variability and directed drive}

We hypothesized that telencephalic circuits are functionally distinguished from other zebrafish brain divisions by a profile dominated by high FC variability and directed output drive. To test this hypothesis, we quantified five regional measures: functional connectivity variability (FCV), functional connectivity strength (FCS), FC partner reconfiguration, $\mathrm{TE}_{\mathrm{net}}$ (outgoing minus incoming directed information flow), and neighborhood $\mathrm{TE}_{\mathrm{net}}$ (the mean directed-drive profile of significantly FC-connected neighbors; definitions in Methods; Fig.~\ref{fig:fc_analysis}).

Hierarchical clustering of these five measures revealed a clear division-level organization of FC-dynamics profiles (Fig.~\ref{fig:fc_analysis}A). Telencephalic regions clustered with elevated FCV, high FC partner reconfiguration and stronger positive $\mathrm{TE}_{\mathrm{net}}$, whereas hindbrain regions tended to show lower FCV and FC partner reconfiguration but stronger mean FCS. Thus, telencephalic circuits were distinguished less by uniformly stronger coupling and more by variable, reconfigurable, and directionally biased interactions.

Division-level comparisons using recording-by-region observations confirmed this pattern (Kruskal--Wallis, all $p < 0.01$; Fig.~\ref{fig:fc_analysis}C--G). Telencephalic regions displayed higher zFCV ($0.70 \pm 0.15$) than diencephalic ($0.13 \pm 0.08$), mesencephalic ($-0.17 \pm 0.09$), and hindbrain regions ($-0.24 \pm 0.06$; all Tel comparisons $p_{\mathrm{Holm}} < 0.001$). In contrast, telencephalic regions showed lower zFCS ($-0.88 \pm 0.10$) than diencephalic ($0.38 \pm 0.08$) and hindbrain regions ($0.26 \pm 0.06$; both $p_{\mathrm{Holm}} < 0.001$). Thus, telencephalic circuits were distinguished by a high-variability, low-mean-coupling FC profile, rather than by uniformly stronger functional coupling.

FC partner reconfiguration further separated divisions. Telencephalic regions showed high profile distance ($0.66 \pm 0.10$), comparable to mesencephalic regions ($0.70 \pm 0.10$), and significantly higher than hindbrain regions ($-0.34 \pm 0.06$; $p_{\mathrm{Holm}} < 0.001$). This indicates that telencephalic FC profiles varied in partner structure across time rather than only changing in overall amplitude.

Directed interaction measures showed a related telencephalic bias. Telencephalic regions had higher $\mathrm{TE}_{\mathrm{net}}$ ($0.53 \pm 0.13$) than diencephalic regions ($-0.62 \pm 0.13$; $p_{\mathrm{Holm}} < 0.001$) and hindbrain regions ($0.09 \pm 0.07$; $p_{\mathrm{Holm}} = 0.019$). Neighborhood $\mathrm{TE}_{\mathrm{net}}$ was also elevated in telencephalon ($0.43 \pm 0.16$) relative to diencephalon ($0.10 \pm 0.10$; $p_{\mathrm{Holm}} = 0.046$), mesencephalon ($-1.15 \pm 0.34$; $p_{\mathrm{Holm}} < 0.001$), and hindbrain ($-0.02 \pm 0.07$; $p_{\mathrm{Holm}} = 0.009$). Together, these directed-interaction measures indicate that telencephalic regions are not only functionally variable, but also embedded in a directed-interaction regime biased toward outgoing information flow and coupled to neighbors with similarly elevated directed-drive profiles.

A regional directed interaction network reconstructed from FC and transfer-entropy data further supported this interpretation (Fig.~\ref{fig:fc_analysis}B). In this network, high-FCV telencephalic nodes participated in FC-supported outgoing transfer paths to other divisions. Together, these results indicate that zebrafish telencephalic circuits occupy a distinctive functional regime characterized by high temporal FC variability, reconfigurable FC partner profiles, and a directed-interaction profile biased toward outgoing information flow.

\begin{figure*}[!ht]
\includegraphics[width=\linewidth]{figure/figure12_final.png}
\caption{
\textbf{Structural directional embedding and second-order motifs distinguish telencephalic circuits.}
(A) Hierarchical clustering of region-wise structural features.
Rows show post-synaptic directional coreness asymmetry ($\mathrm{DCA}_{\mathrm{post}}$), pre-synaptic directional coreness asymmetry ($\mathrm{DCA}_{\mathrm{pre}}$), modularity Q, log-transformed out/in ratio, and output--output motif fraction, each z-scored across regions.
Columns indicate anatomical regions ordered by Ward hierarchical clustering.
Regions are grouped into telencephalic (Tel), diencephalic (Di), mesencephalic (Mes), and hindbrain (Hind), as indicated by the top annotation bar.
(B) Directed region-level structural connectivity network matched to the regions shown in panel A.
Nodes represent anatomical regions and are arranged by anatomical division along the vertical axis and by mean $\mathrm{DCA}_{\mathrm{post}}$ along the horizontal axis, with high $\mathrm{DCA}_{\mathrm{post}}$ regions on the left and low $\mathrm{DCA}_{\mathrm{post}}$ regions on the right.
Directed edges represent region-level structural connectivity aggregated from cell-level synaptic connections across subjects.
Node colors indicate anatomical divisions.
(C) Representative directed structural connectivity patterns illustrating directional coreness asymmetry between pairs of brain regions.
Nodes represent neurons arranged in a circular layout.
Node colors encode neuronal DCA values, with red indicating larger DCA values and blue indicating smaller DCA values.
Directed edges indicate synaptic connections, with orange edges projecting toward the lower region and blue edges projecting toward the upper region.
Examples include pallium to intermediate reticular formation (P--imRF), pallium to medulla oblongata strip 4 (P--MOS4), raphe to posterior reticular formation (Ra--pRF), and raphe to intermediate reticular formation (Ra--imRF).
(D--H) Division-level distributions of $\mathrm{DCA}_{\mathrm{post}}$ (D), $\mathrm{DCA}_{\mathrm{pre}}$ (E), modularity Q (F), log-transformed out/in ratio (G), and output--output motif fraction (H).
Box plots show median and interquartile range; points indicate subject-region observations.
Group differences were assessed using Kruskal--Wallis tests followed by Holm-corrected two-sided Mann--Whitney U tests.
Kruskal--Wallis results were: $\mathrm{DCA}_{\mathrm{post}}$, $H = 135.359$, $p = 3.78\times 10^{-29}$, $\eta^2_H = 0.315$; $\mathrm{DCA}_{\mathrm{pre}}$, $H = 80.039$, $p = 3.01\times 10^{-17}$, $\eta^2_H = 0.184$; modularity Q, $H = 143.833$, $p = 5.63\times 10^{-31}$, $\eta^2_H = 0.435$; log-transformed out/in ratio, $H = 11.973$, $p = 7.48\times 10^{-3}$, $\eta^2_H = 0.020$; and output--output motif fraction, $H = 74.998$, $p = 3.63\times 10^{-16}$, $\eta^2_H = 0.250$.
Statistical significance is indicated as *$p < 0.05$, **$p < 0.01$, and ***$p < 0.001$.
}
\label{fig:sc_analysis}
\end{figure*}

\subsection*{Telencephalic regions exhibit output-biased second-order cascade motifs}

We next asked which structural-connectivity (SC) features were associated with regional FCV. We hypothesized that FCV reflects cascaded directional embedding, in which regional outputs preferentially target downstream populations that are themselves positioned to broadcast signals onward. To test this hypothesis, we reconstructed a synapse-informed directed SC network by combining neuronal soma coordinates from calcium imaging \citep{chen2018brain} with synapse-like endpoints from a cellular-resolution anatomical dataset \citep{kunst2019cellular} (Methods).

From this SC network, we constructed a region-wise structural feature profile consisting of post-synaptic directional coreness asymmetry ($\mathrm{DCA}_{\mathrm{post}}$), pre-synaptic directional coreness asymmetry ($\mathrm{DCA}_{\mathrm{pre}}$), modularity Q, log-transformed out/in ratio, and output--output motif fraction (OO-fraction; definitions in Methods). These measures jointly capture directional embedding, modular segregation, net output/input bias, and second-order synaptic cascade structure.

$\mathrm{DCA}_{\mathrm{pre}}$ and $\mathrm{DCA}_{\mathrm{post}}$ quantify directional embedding on opposite sides of a region's connectivity.
$\mathrm{DCA}_{\mathrm{pre}}$ describes the input side of a region: it asks whether neurons projecting into that region originate from output-biased or input-biased local circuits.
A higher $\mathrm{DCA}_{\mathrm{pre}}$ therefore indicates that a region receives input from upstream neurons that are themselves positioned to broadcast signals onward.
$\mathrm{DCA}_{\mathrm{post}}$ describes the output side of a region: it asks whether neurons targeted by that region's outputs are output-biased or input-biased within their own local circuits.
A higher $\mathrm{DCA}_{\mathrm{post}}$ therefore indicates that a region sends projections to downstream neurons positioned to propagate signals further, consistent with outward, cascade-like embedding.
OO-fraction provides a complementary second-order motif measure by quantifying how often a region's outgoing projections target neurons that themselves have outgoing structural partners.

Comparing outgoing projection targets across paired regions---P to imRF or MOS4, and Ra to pRF or imRF---revealed qualitative differences in synaptic target embedding (Fig.~\ref{fig:sc_analysis}C).
P outgoing projections landed predominantly on target neurons with higher local DCA values, yielding higher regional $\mathrm{DCA}_{\mathrm{post}}$, whereas Ra projections targeted neurons with lower local DCA values, yielding lower regional $\mathrm{DCA}_{\mathrm{post}}$.
Thus, regional differences in structural embedding depended not only on whether a region sent outputs, but also on the cascade embedding of the neurons those outputs targeted.

Hierarchical clustering of the five structural features revealed division-specific organization of structural connectivity (Fig.~\ref{fig:sc_analysis}A).
Telencephalic regions were characterized by elevated $\mathrm{DCA}_{\mathrm{post}}$, high OO-fraction, and relatively high modularity Q.
This profile suggests that telencephalic circuits are not simply distinguished by output degree alone, but by a structural arrangement in which outputs preferentially reach downstream neurons embedded in further outgoing cascades.

Division-level comparisons confirmed significant differences across all five structural measures (Kruskal--Wallis, all $p < 0.01$; Fig.~\ref{fig:sc_analysis}D--H).

$\mathrm{DCA}_{\mathrm{post}}$ showed a clear telencephalic enrichment. Telencephalic regions had higher $\mathrm{DCA}_{\mathrm{post}}$ ($0.008 \pm 0.005$) than diencephalic ($-0.011 \pm 0.001$), mesencephalic ($-0.025 \pm 0.003$), and hindbrain regions ($-0.012 \pm 0.001$; all $p_{\mathrm{Holm}} < 0.001$). This indicates preferential targeting of output-biased downstream populations.

$\mathrm{DCA}_{\mathrm{pre}}$ also differed across divisions, but did not show the same telencephalic enrichment as $\mathrm{DCA}_{\mathrm{post}}$. Hindbrain regions showed a higher $\mathrm{DCA}_{\mathrm{pre}}$ distribution than telencephalic and mesencephalic regions (both $p_{\mathrm{Holm}} < 0.001$), whereas telencephalic regions did not differ significantly from diencephalic or mesencephalic regions. This indicates that input-side embedding varied across divisions, but was not the primary telencephalic structural signature.

Modularity Q also differed across divisions. Telencephalic regions showed higher modularity Q ($0.696 \pm 0.010$) than hindbrain regions ($0.523 \pm 0.009$; $p_{\mathrm{Holm}} < 0.001$), whereas mesencephalic regions showed the highest modularity overall ($0.736 \pm 0.014$).

The log-transformed out/in ratio showed a weaker but significant division-level effect. Telencephalic regions showed a distributional shift in log$_{10}$(out/in) relative to diencephalic and hindbrain regions ($p_{\mathrm{Holm}} < 0.05$ and $p_{\mathrm{Holm}} < 0.05$, respectively), but did not differ from mesencephalic regions. Because this effect was weaker and less monotonic than $\mathrm{DCA}_{\mathrm{post}}$ or OO-fraction, first-order output/input imbalance alone did not explain the full structural pattern.

OO-fraction showed a strong division-level difference. Telencephalic regions exhibited higher OO-fraction ($0.264 \pm 0.009$) than diencephalic ($0.230 \pm 0.009$; $p_{\mathrm{Holm}} = 0.0052$), mesencephalic ($0.164 \pm 0.008$; $p_{\mathrm{Holm}} < 0.001$), and hindbrain regions ($0.185 \pm 0.003$; $p_{\mathrm{Holm}} < 0.001$). 

A region-level directed SC network provided a complementary view of this organization (Fig.~\ref{fig:sc_analysis}B).
Regions were arranged by anatomical division and ordered along the horizontal axis by mean $\mathrm{DCA}_{\mathrm{post}}$, with higher $\mathrm{DCA}_{\mathrm{post}}$ regions positioned toward the left.
In this embedding, telencephalic regions were shifted toward the high $\mathrm{DCA}_{\mathrm{post}}$ side of the network, indicating stronger downstream targeting of output-biased populations within the reconstructed structural embedding. 



\begin{figure*}[!ht]
\includegraphics[width=\linewidth]{figure/figure_sc_fc_final_overview.png}
\caption{
\textbf{Directional structural embedding and second-order output motifs are associated with regional FC variability.}
(A) Pairwise correlations between regional functional features and structural features.
Functional features are the five measures used in Fig.~\ref{fig:fc_analysis}: functional connectivity variability (FCV), functional connectivity strength (FCS), FC profile correlation-distance, net transfer entropy ($\mathrm{TE}_{\mathrm{net}}$), and neighborhood $\mathrm{TE}_{\mathrm{net}}$.
Structural features are the five measures used in Fig.~\ref{fig:sc_analysis}: postsynaptic directional coreness asymmetry ($\mathrm{DCA}_{\mathrm{post}}$), presynaptic directional coreness asymmetry ($\mathrm{DCA}_{\mathrm{pre}}$), modularity Q, log-transformed out/in ratio, and output--output motif fraction.
Color indicates Pearson correlation coefficient ($r$).
Black outlines mark correlations that remained significant after Benjamini--Hochberg FDR correction across all 25 feature pairs.
(B) First latent variable (LV1) from partial least squares (PLS) analysis linking the functional and structural feature spaces.
Each point represents a brain region and is colored by major brain division: telencephalic (Tel), diencephalic (Di), mesencephalic (Mes), and hindbrain (Hind).
(C) Functional feature weights on PLS LV1.
Bars indicate the mean leave-one-region-out LV1 weights, points indicate individual leave-one-region-out estimates, and diamonds with whiskers indicate the mean and 95\% confidence interval of the leave-one-region-out estimates.
(D) Structural feature weights on PLS LV1, shown with the same leave-one-region-out stability analysis as in (C).
(E) Cross-validated prediction of FCV from structural features using linear regression.
Predicted FCV values were obtained using fivefold cross-validation, and the dashed line indicates identity.
(F) Standardized regression coefficients for predicting FCV from structural features, shown with the same leave-one-region-out stability analysis as in (C).
}
\label{fig:dca_vs_fcv}
\end{figure*}

\subsection*{Outgoing cascade structure is associated with regional FC variability}

Having established that telencephalic circuits exhibit elevated FC variability and directed interaction profiles, we next asked which features of the directed structural connectome were associated with regional FCV variation. We quantified pairwise associations between the five regional functional features used in Fig.~\ref{fig:fc_analysis}---FCV, FCS, FC partner reconfiguration, $\mathrm{TE}_{\mathrm{net}}$, and neighborhood $\mathrm{TE}_{\mathrm{net}}$---and the five structural features used in Fig.~\ref{fig:sc_analysis}: $\mathrm{DCA}_{\mathrm{post}}$, $\mathrm{DCA}_{\mathrm{pre}}$, modularity Q, log-transformed out/in ratio, and output--output motif fraction.

Across all 25 structure--function pairs, FCV showed the strongest positive associations with $\mathrm{DCA}_{\mathrm{post}}$ and output--output motif fraction (Fig.~\ref{fig:dca_vs_fcv}A). Regions with higher $\mathrm{DCA}_{\mathrm{post}}$ had higher FCV (Pearson's $r = 0.577$, $q_{\mathrm{FDR}} < 0.001$), indicating that regions whose outputs target more output-biased downstream populations show greater temporal reconfiguration of functional coupling. FCV was also strongly associated with output--output motif fraction ($r = 0.560$, $q_{\mathrm{FDR}} < 0.001$), suggesting that second-order output cascade motifs capture a closely related structural signature. In contrast, $\mathrm{DCA}_{\mathrm{pre}}$ was not positively associated with FCV ($r = -0.170$, $q_{\mathrm{FDR}} = 0.368$), indicating that the FCV-linked structural signal was more prominent on the outgoing target side than on the incoming source side.


To examine multivariate structure--function coupling, we performed PLS analysis between the full functional and structural feature spaces. The first latent variable revealed a strong shared structure--function axis ($r = 0.685$, $p < 0.001$; Fig.~\ref{fig:dca_vs_fcv}B). The functional side of this axis was dominated by high FCV, FC partner reconfiguration, and lower FCS (Fig.~\ref{fig:dca_vs_fcv}C). The structural side was dominated by high $\mathrm{DCA}_{\mathrm{post}}$ , high output--output motif fraction, and high log-transformed out/in ratio, with a negative contribution from $\mathrm{DCA}_{\mathrm{pre}}$ (Fig.~\ref{fig:dca_vs_fcv}D). Thus, regions with more variable functional coupling occupied a structural regime characterized by outgoing cascade embedding and enriched second-order output motifs.


We then tested whether structural features could predict FCV directly. A fivefold cross-validated linear model showed modest prediction of regional FCV from structural features (observed--predicted FCV: $r = 0.441$, $p < 0.01$; $R^2_{\mathrm{CV}} = 0.134$; Fig.~\ref{fig:dca_vs_fcv}E). Standardized coefficients again identified $\mathrm{DCA}_{\mathrm{post}}$ and output--output motif fraction as the strongest positive predictors of FCV ($\beta = 0.311$ and $\beta = 0.300$, respectively), whereas $\mathrm{DCA}_{\mathrm{pre}}$ contributed weakly in the negative direction ($\beta = -0.061$; Fig.~\ref{fig:dca_vs_fcv}F). Modularity Q and log-transformed out/in ratio contributed only weakly in this direct FCV prediction model.

Together, these analyses indicate that regional FCV is associated most strongly with an output-oriented structural axis at the region level, rather than with generic modularity or first-order output/input imbalance alone. In particular, $\mathrm{DCA}_{\mathrm{post}}$ and OO-fraction point to the same organization: high-FCV regions preferentially project to downstream targets that are themselves embedded in output-biased cascade motifs. Thus, flexible functional coupling is linked to second-order outgoing cascade structure.



\begin{figure*}[!ht]
\includegraphics[width=\linewidth]{figure/figure13_final.png}
\caption{
\textbf{Layer-wise asymmetry and direction-specific rewiring link directed structure to functional connectivity variability.}
(A) Schematic of the layer linear model.
Four layers with increasing node number are connected by symmetric within-layer coupling and asymmetric between-layer directional coupling. As the asymmetry parameter $\epsilon$ increases, downward coupling is strengthened and upward coupling is weakened. $\epsilon$ was used as an abstract analogue of downstream output-dominated embedding.
(B) Effective energy potential landscapes of the layer linear model.
For Layer 1 and Layer 4, dynamic FC states were estimated from repeated simulations and converted into empirical potentials, $U(z)=-\log P(z)$, where $P(z)$ is the probability distribution of the dynamic FC state across sliding windows.
Potential curves were estimated separately for each $\epsilon$ and summarized across 50 simulation runs; lines show the mean and shaded bands show SEM.
Increasing $\epsilon$ reshaped the potential landscape, indicating that directed layer asymmetry changes the stability and accessibility of dynamic FC states.
(C) Layer-mean functional connectivity variability (FCV) across the asymmetry parameter $\epsilon$.
(D) Whole-brain zebrafish simulation and direction-specific rewiring schematic centered on P/SP neurons.
The base model uses the empirical synapse-informed connectome involving P/SP neurons and their connected partners. NULL-Out rewires outgoing edges from P/SP neurons, whereas NULL-In rewires incoming edges to P/SP neurons. In both null conditions, anatomical region pairs are preserved while cell-level synaptic endpoints are shuffled within the same regions.
(E) Bootstrap distributions of direction-specific rewiring effects on P/SP $\Delta \mathrm{FCV}$ in the zebrafish simulation.
For each condition, FCV was first averaged across all P/SP nodes, and $\Delta \mathrm{FCV}$ was defined as the difference between these P/SP mean zFCV values; bootstrap replicates were generated by resampling P/SP nodes within each condition.
Distributions show NULL-Out minus Base, NULL-In minus Base, and NULL-Out minus NULL-In differences. NULL-Out produced a larger reduction in FCV than NULL-In
($\Delta_{\mathrm{Out-Base}} = -0.558$, 95\% CI = $[-0.747, -0.365]$, $p_{\mathrm{Holm}} < 0.001$;
$\Delta_{\mathrm{In-Base}} = -0.316$, 95\% CI = $[-0.510, -0.123]$, $p_{\mathrm{Holm}} = 0.004$;
$\Delta_{\mathrm{Out-In}} = -0.242$, 95\% CI = $[-0.453, -0.027]$, $p_{\mathrm{Holm}} = 0.029$).
Stars indicate Holm-corrected bootstrap significance: $^{*}p < 0.05$, $^{**}p < 0.01$, $^{***}p < 0.001$.
}
\label{fig:layer-linear-model}
\end{figure*}

\subsection*{Directed cascade simulations support direction-specific model sensitivity of FC variability}

To test whether directional cascade organization can provide mechanistic plausibility for FC variability, we first used a simplified layer-wise linear model. The model consisted of four layers with increasing layer size. Within-layer coupling was symmetric, whereas between-layer coupling was asymmetric: downward coupling was strengthened and upward coupling was weakened as the asymmetry parameter $\epsilon$ increased (Fig.~\ref{fig:layer-linear-model}A). This model was not intended to reproduce the full zebrafish connectome, but to isolate the effect of directional imbalance in a minimal cascade-like architecture. We treated $\epsilon$ as an abstract model analogue of outgoing cascade bias, conceptually related to $\mathrm{DCA}_{\mathrm{post}}$ and to second-order output-to-output motif enrichment. Full equations and representative layer activity traces are provided in SI Appendix, Four-layer Directed Linear Network Model: Full Specification and Fig.~S8B.

We next asked how directional asymmetry reshaped the dynamic FC-state landscape. For each simulation, we defined a layer-wise FC state as the sliding-window mean FC between one layer and all other layers, and converted its empirical distribution into an effective potential. Increasing $\epsilon$ altered this landscape, especially in Layer 1, where FC states became broader and shifted across potential wells (Fig.~\ref{fig:layer-linear-model}B). Thus, stronger outgoing cascade bias made the system less confined to a single dominant FC regime and increased its tendency to sample alternative connectivity states over time.

Consistent with this potential-landscape view, layer-wise FCV increased with increasing $\epsilon$, with the strongest effect observed in the first layer (Fig.~\ref{fig:layer-linear-model}C). These results indicate that asymmetric inter-layer coupling can increase temporal variability in functional connectivity by broadening the accessible range of transient FC configurations.

We next tested whether this direction-dependent effect was also present in a zebrafish whole-brain network model based on the synapse-informed directed SC network. The simulation focused on pallium/subpallium (P/SP) neurons, and two direction-specific null models were constructed by rewiring cell-level synaptic endpoints while preserving anatomical region pairs (Fig.~\ref{fig:layer-linear-model}D). In NULL-Out, outgoing edges from P/SP neurons were rewired; in NULL-In, incoming edges to P/SP neurons were rewired. Thus, both null models preserved coarse anatomical connectivity while selectively disrupting the cell-level organization of either P/SP outgoing or incoming projections. Simulation equations, weight optimization, empirical validation, and null-control examples are provided in SI Appendix, Large-scale Neuronal Simulations: Model Equations; Large-scale Network Simulation: Weight Optimization and Empirical Validation; and Figs.~S6--S8.

Direction-specific rewiring produced a clear effect on P/SP FCV (Fig.~\ref{fig:layer-linear-model}E). Both outgoing and incoming P/SP rewiring reduced FCV relative to the base network, but the outgoing perturbation produced the larger reduction (NULL-Out: $\Delta \mathrm{FCV} = -0.558$, $p_{\mathrm{Holm}} < 0.001$; NULL-In: $\Delta \mathrm{FCV} = -0.316$, $p_{\mathrm{Holm}} < 0.01$; NULL-Out vs. NULL-In: $\Delta \mathrm{FCV} = -0.242$, $p_{\mathrm{Holm}} < 0.05$). Thus, in this model, P/SP FCV was more sensitive to outgoing than incoming rewiring.

Together, the minimal layer model and the zebrafish whole-brain simulation provide mechanistic plausibility for an outgoing cascade mechanism for FCV. Increasing $\epsilon$, an abstract analogue of second-order outgoing motif bias, broadened the accessible FC-state landscape and increased FCV in simplified dynamics, whereas disrupting outgoing P/SP projections preferentially reduced FCV in the zebrafish network model. This agreement is consistent with outgoing projection organization as a candidate structural contributor to temporally flexible functional connectivity.


\begin{figure*}[!ht]
\includegraphics[width=\linewidth]{figure/figure14_15_celegans_drosophila_AFHI_combined.png}
\caption{
\textbf{Directed structural embedding and FC variability in \textit{C.~elegans} and adult \textit{Drosophila melanogaster}.}
Top row, \textit{C.~elegans}; bottom row, \textit{Drosophila}.
(A) Directed chemical synapse network of the final analyzed \textit{C.~elegans} neurons, with nodes colored by anatomical class and arranged horizontally by $\mathrm{DCA}_{\mathrm{post}}$.
(B) Neuron-level $\mathrm{DCA}_{\mathrm{post}}$ versus recording-averaged z-scored FCV in \textit{C.~elegans} (Pearson $r = 0.170$, $p = 0.071$).
(C) Recording-level FCV distributions across \textit{C.~elegans} anatomical classes.
(D,E) Neuron-level $\mathrm{DCA}_{\mathrm{post}}$ and $\mathrm{DCA}_{\mathrm{pre}}$ distributions across \textit{C.~elegans} classes.
(F) Directed FlyWire structural network of matched \textit{Drosophila} Ito-region nodes, colored by anatomical division and arranged horizontally by $\mathrm{DCA}_{\mathrm{pre}}$.
(G) Region-level $\mathrm{DCA}_{\mathrm{pre}}$ versus recording-averaged z-scored FCV in \textit{Drosophila} (Pearson $r = 0.580$, $p = 1.84\times 10^{-4}$).
(H) Recording-level FCV distributions across \textit{Drosophila} anatomical divisions.
(I,J) Region-level $\mathrm{DCA}_{\mathrm{post}}$ and $\mathrm{DCA}_{\mathrm{pre}}$ distributions across \textit{Drosophila} anatomical divisions.
For group panels, points show individual recording-level FCV observations or region-level SC observations, and boxes show median and interquartile range.
For zFCV panels (C,H), global group differences were assessed using Kruskal--Wallis tests, and pairwise significance bars indicate Holm-corrected two-sided Mann--Whitney U tests; panel H shows mushroom body-centered comparisons with selected olfactory neuropil comparisons for readability.
For DCA panels (D,E,I,J), global group differences were assessed using group-label permutation tests, and pairwise significance bars indicate Holm-corrected permutation tests.
Global test results were:
\textit{C.~elegans} zFCV, Kruskal--Wallis $H = 45.040$, $\eta^2_H = 0.032$, $p = 3.90\times 10^{-9}$;
\textit{C.~elegans} $\mathrm{DCA}_{\mathrm{post}}$, permutation $p = 1.50\times 10^{-4}$;
\textit{C.~elegans} $\mathrm{DCA}_{\mathrm{pre}}$, permutation $p = 0.0668$;
\textit{Drosophila} zFCV, Kruskal--Wallis $H = 306.842$, $\eta^2_H = 0.229$, $p = 3.38\times 10^{-64}$;
\textit{Drosophila} $\mathrm{DCA}_{\mathrm{post}}$, permutation $p = 5.50\times 10^{-4}$;
and \textit{Drosophila} $\mathrm{DCA}_{\mathrm{pre}}$, permutation $p = 0.0105$.
Statistical significance is indicated as *$p < 0.05$, **$p < 0.01$, and ***$p < 0.001$.
}
\label{fig:celegans_drosophila}
\end{figure*}


\subsection*{Comparative analyses reveal species- and scale-dependent directed embedding signatures of FCV}

We next asked whether directed structural embedding is associated with functional variability in nervous systems with very different organizational scales. We focused on the fully identified \textit{C.~elegans} connectome \citep{Cook2019,AtanasKim2023} and the adult \textit{Drosophila melanogaster} brain \citep{Dorkenwald2024FlyWire,Turner2021Branson999} because both systems provide matched structural connectivity and functional activity data at biologically interpretable scales. In both systems, we used DCA to quantify output-biased versus input-biased directed embedding, and FCV was computed from sliding-window correlations of calcium activity (Methods).

In \textit{C.~elegans}, we restricted the structure--function analysis to neurons with sufficient spontaneous recording coverage, yielding 120 neurons for the final neuron-level comparison. The directed chemical synapse network was visualized with neurons colored by anatomical class and arranged horizontally by $\mathrm{DCA}_{\mathrm{post}}$ (Fig.~\ref{fig:celegans_drosophila}A). Across neurons, $\mathrm{DCA}_{\mathrm{post}}$ showed a weak positive association with recording-averaged zFCV (Pearson $r = 0.170$, $p = 0.071$; Fig.~\ref{fig:celegans_drosophila}B). Although this neuron-level trend did not reach conventional significance, it was consistent in direction with the idea that neurons whose outputs target more output-biased downstream partners tend to show greater temporal variability in functional coupling.

Anatomical class analyses revealed stronger group-level organization. Recording-level zFCV differed significantly across \textit{C.~elegans} classes (Fig.~\ref{fig:celegans_drosophila}C). Olfactory/chemosensory neurons showed elevated zFCV ($0.286 \pm 0.066$), significantly higher than sensory neurons and interneurons (both $p_{\mathrm{Holm}} < 0.001$). Structural embedding also differed across classes. $\mathrm{DCA}_{\mathrm{post}}$ showed significant class-level differences (Fig.~\ref{fig:celegans_drosophila}D), with olfactory/chemosensory neurons showing positive $\mathrm{DCA}_{\mathrm{post}}$ ($0.014 \pm 0.006$) and higher values than sensory, motor/command, and state-modulatory neurons (all $p_{\mathrm{Holm}} < 0.05$). In contrast, $\mathrm{DCA}_{\mathrm{pre}}$ showed weaker class-level separation (Fig.~\ref{fig:celegans_drosophila}E), with no significant pairwise class separation. Thus, in \textit{C.~elegans}, olfactory/chemosensory neurons showed both elevated FCV and elevated $\mathrm{DCA}_{\mathrm{post}}$, indicating that their high functional variability is accompanied by stronger postsynaptic downstream embedding.

We then asked whether a related association was present in the adult fly brain. Because neuron-by-neuron identity matching between FlyWire structural data and Branson functional calcium recordings is not available, the analysis was performed at the level of matched side-aware Ito anatomical regions. The directed FlyWire structural network was visualized with nodes colored by anatomical division and arranged horizontally by $\mathrm{DCA}_{\mathrm{pre}}$ (Fig.~\ref{fig:celegans_drosophila}F), highlighting the input-side profile observed in high-FCV mushroom body regions.


Across matched \textit{Drosophila} regions, $\mathrm{DCA}_{\mathrm{pre}}$ was positively associated with recording-averaged FCV (Pearson $r = 0.580$, $p < 0.001$; Fig.~\ref{fig:celegans_drosophila}G), whereas $\mathrm{DCA}_{\mathrm{post}}$ was not significantly associated with FCV ($r = -0.260$, $p = 0.120$). However, these region-level DCA--FCV relationships should be interpreted cautiously because they were strongly influenced by the distinct values of mushroom body regions. We therefore do not emphasize the fly result as evidence for a broad continuous DCA--FCV relationship across all regions. Instead, the more relevant observation is that high-FCV mushroom body regions showed a direction-specific structural profile, with high $\mathrm{DCA}_{\mathrm{pre}}$ and low $\mathrm{DCA}_{\mathrm{post}}$.


Recording-level FCV differed strongly across \textit{Drosophila} anatomical divisions (Fig.~\ref{fig:celegans_drosophila}H). Mushroom body regions showed the highest FCV ($1.063 \pm 0.138$), significantly exceeding optic neuropil, central complex, protocerebrum, and premotor interface regions (all $p_{\mathrm{Holm}} < 0.001$), as well as olfactory neuropil ($p_{\mathrm{Holm}} < 0.05$). Structural group analyses showed opposite direction-specific tendencies (Fig.~\ref{fig:celegans_drosophila}I,J). The global division effect for $\mathrm{DCA}_{\mathrm{post}}$ was not significant ($p_{\mathrm{perm}} = 0.125$), but mushroom body regions showed low mean $\mathrm{DCA}_{\mathrm{post}}$ ($-0.091 \pm 0.034$). By contrast, $\mathrm{DCA}_{\mathrm{pre}}$ showed a significant division-level effect ($p_{\mathrm{perm}} = 0.0385$), with mushroom body regions showing high $\mathrm{DCA}_{\mathrm{pre}}$, although no pairwise division contrast survived Holm correction. Thus, in \textit{Drosophila}, the high-FCV mushroom body profile was accompanied by high input-side and low output-side directed embedding.

Together with the zebrafish results, these analyses suggest that high-FCV groups or regions can be accompanied by directed structural embedding, but the relevant side of the directed architecture differs across species and analytical scale. In zebrafish, high FCV in telencephalic regions was most strongly linked to outgoing downstream embedding, captured by elevated $\mathrm{DCA}_{\mathrm{post}}$ and output--output motif fraction. In \textit{C.~elegans}, the clearest group-level signature similarly involved elevated $\mathrm{DCA}_{\mathrm{post}}$ in high-FCV olfactory/chemosensory neurons. In \textit{Drosophila}, by contrast, high-FCV mushroom body regions showed an opposite directional profile, with high $\mathrm{DCA}_{\mathrm{pre}}$ and low $\mathrm{DCA}_{\mathrm{post}}$. Overall, the comparative evidence supports a species- and scale-dependent pattern rather than a single conserved cross-species rule.



\begin{figure}[!ht]
\includegraphics[width=\linewidth]{figure/figure7_by_species_with_figure5_abc.png}
\caption{
\textbf{Condition-associated FCV relationships and functional-group structure--function signatures across species.}
(A--C) Relationship between spontaneous FCV and condition-associated FCV summaries.
In \textit{C.~elegans}, spontaneous FCV is compared with heat-epoch FCV (A).
In zebrafish, spontaneous FCV is compared with the across-condition standard deviation of OMR-associated FCV (B) and mean OMR-associated FCV (C).Points indicate neurons or anatomical regions depending on the available matched scale; point colors indicate functional groups.
(D--F) Functional-group signatures across species for FCV (D), $\mathrm{DCA}_{\mathrm{post}}$ (E), and $\mathrm{DCA}_{\mathrm{pre}}$ (F).
Columns show \textit{C.~elegans}, \textit{Drosophila}, and zebrafish.
Within each species and metric, values were centered and scaled to emphasize relative functional-group enrichment.
Violin plots show the distribution of individual recording-, neuron-, region-, or subject-region observations depending on the available matched scale; large markers show group means, connected by lines to visualize group-level trends.
Marker shapes indicate species: circles, \textit{C.~elegans}; squares, \textit{Drosophila}; triangles, zebrafish.
Functional groups are shown along the x axis, with olfactory/chemosensory and associative/state-related groups highlighted by color.
For each species and metric, functional-group enrichment was tested as the absolute group-versus-rest standardized mean difference after permuting functional-group labels while preserving group sizes; $p$ values were Holm-corrected across groups within each species--metric comparison.
Stars indicate Holm-corrected enrichment:
*$p < 0.05$, **$p < 0.01$, and ***$p < 0.001$.
}
\label{fig:structure_function_synthesis}
\end{figure}


\subsection*{Functional-group synthesis reveals stimulus-linked and species-specific cascade signatures of FCV}

First, we tested whether spontaneous FCV was related to FCV measured under stimulus conditions (Fig.~\ref{fig:structure_function_synthesis}A--C). This analysis was limited to \textit{C.~elegans} and zebrafish because the available \textit{Drosophila} calcium imaging dataset consisted of intrinsic recordings acquired without sensory stimulation \citep{Turner2021Branson999}. In \textit{C.~elegans}, spontaneous FCV was strongly related to heat-epoch FCV ($r = 0.84$, $p < 0.001$; Fig.~\ref{fig:structure_function_synthesis}A). In zebrafish, spontaneous FCV was positively related to the across-condition standard deviation of OMR-associated FCV ($r = 0.36$, $p < 0.05$; Fig.~\ref{fig:structure_function_synthesis}B) and more strongly related to mean OMR-associated FCV ($r = 0.62$, $p < 0.001$; Fig.~\ref{fig:structure_function_synthesis}C). Thus, regions with higher spontaneous FCV also tended to show higher stimulus-epoch FCV, with a weaker but significant association for across-condition FCV modulation.


We next summarized FCV and directed structural embedding across five broad functional groups in each species: olfactory/chemosensory, sensory-related, sensorimotor, integrative, and associative/state-related groups (Fig.~\ref{fig:structure_function_synthesis}D--F; see Methods for classification criteria). In these panels, rows show FCV (Fig.~\ref{fig:structure_function_synthesis}D), $\mathrm{DCA}_{\mathrm{post}}$ (Fig.~\ref{fig:structure_function_synthesis}E), and $\mathrm{DCA}_{\mathrm{pre}}$ (Fig.~\ref{fig:structure_function_synthesis}F), while columns show \textit{C.~elegans}, \textit{Drosophila}, and zebrafish. Olfactory/chemosensory and associative/state-related groups showed prominent FCV enrichment, but their associated directed structural profiles differed across species and analysis scales.

In \textit{C.~elegans}, olfactory/chemosensory neurons showed elevated FCV (relative $z=0.72$, $p_{\mathrm{Holm}} < 0.01$; Fig.~\ref{fig:structure_function_synthesis}D) together with elevated $\mathrm{DCA}_{\mathrm{post}}$ (relative $z=0.74$, $p_{\mathrm{Holm}} < 0.05$; Fig.~\ref{fig:structure_function_synthesis}E), whereas $\mathrm{DCA}_{\mathrm{pre}}$ was not similarly elevated (Fig.~\ref{fig:structure_function_synthesis}F). Thus, at the identified-neuron scale, the high-FCV chemosensory group was accompanied primarily by an output-side directed embedding tendency.

In zebrafish, associative/state-related regions showed FCV enrichment (relative $z=0.56$, $p_{\mathrm{Holm}} <0.05$; Fig.~\ref{fig:structure_function_synthesis}D) together with elevated $\mathrm{DCA}_{\mathrm{post}}$ (relative $z=1.06$, $p_{\mathrm{Holm}} <0.01$; Fig.~\ref{fig:structure_function_synthesis}E) and reduced $\mathrm{DCA}_{\mathrm{pre}}$ (relative $z=-0.45$, $p_{\mathrm{Holm}} < 0.01$; Fig.~\ref{fig:structure_function_synthesis}F), indicating an output-side directed embedding profile. Olfactory/chemosensory regions showed a related pattern, with elevated FCV and $\mathrm{DCA}_{\mathrm{post}}$ but no significant $\mathrm{DCA}_{\mathrm{pre}}$ enrichment.

In contrast, \textit{Drosophila} associative/state-related regions showed elevated FCV (relative $z=0.73$, $p_{\mathrm{Holm}}=0.005$; Fig.~\ref{fig:structure_function_synthesis}D) together with reduced $\mathrm{DCA}_{\mathrm{post}}$ (relative $z=-3.28$, $p_{\mathrm{Holm}}=0.015$; Fig.~\ref{fig:structure_function_synthesis}E). They also showed a positive $\mathrm{DCA}_{\mathrm{pre}}$ tendency (relative $z=1.73$; uncorrected permutation $p=0.061$), although this did not survive Holm correction ($p_{\mathrm{Holm}}=0.255$), likely reflecting the small number of matched associative/state-related SC nodes in the fly dataset. Thus, the fly associative/state-related group showed a high-FCV profile with a clear low-$\mathrm{DCA}_{\mathrm{post}}$ signature and a suggestive, but not corrected-significant, high-$\mathrm{DCA}_{\mathrm{pre}}$ tendency. Olfactory/chemosensory regions also showed elevated FCV, but their DCA profile was less statistically resolved: $\mathrm{DCA}_{\mathrm{post}}$ was positive but nonsignificant (relative $z=0.93$, uncorrected $p=0.126$), and $\mathrm{DCA}_{\mathrm{pre}}$ showed only a marginal uncorrected tendency (relative $z=1.38$, uncorrected $p=0.051$; $p_{\mathrm{Holm}}=0.255$).


Together, these results suggest two complementary conclusions. First, spontaneous FCV is associated with stimulus-epoch FCV measures, consistent with greater stimulus-linked modulation in high-FCV circuits. Second, high-FCV functional groups are associated with directed cascade-like structural organization, but not with a single invariant DCA direction. The structural route to high FCV appears species- and circuit-dependent: \textit{C.~elegans} olfactory/chemosensory neurons and zebrafish associative/state-related regions align more strongly with output-side $\mathrm{DCA}_{\mathrm{post}}$, whereas \textit{Drosophila} associative/state-related regions align more strongly with low $\mathrm{DCA}_{\mathrm{post}}$ and high $\mathrm{DCA}_{\mathrm{pre}}$.

\section*{Discussion}

The main implication of this study is that functional connectivity variability (FCV) is associated with a second-order structural organization that may support the reconfiguration of circuit interactions over time. In larval zebrafish, high-FCV regions were enriched in the telencephalon and were accompanied by higher net directed-interaction measures, elevated downstream output-biased embedding captured by $\mathrm{DCA}_{\mathrm{post}}$, and enriched output-to-output motif fraction (Figs.~\ref{fig:fc_analysis}--\ref{fig:dca_vs_fcv}). Thus, high-FCV telencephalic circuits were not distinguished simply by having more outgoing projections, but by where those outputs were embedded: their projections preferentially reached downstream populations that were themselves positioned to propagate activity further.

This structural profile is especially relevant for pallial circuits. The zebrafish pallium supports goal-directed behavior, spatial representation, context-dependent adaptation, and associative memory-related synaptic remodeling \citep{tognoli2014metastable,Yang2024SpatialCode,Dempsey2022ZebrafishMemorySynapses}. Anatomically, zebrafish pallial zones also show compartmentalized pallial and extrapallial connections, with region-specific efferent pathways to downstream targets \citep{Yanez2022ZebrafishPalliumHodology}. The second-order cascade topology identified here may therefore provide a candidate structural principle for reconfigurable associative dynamics: not simply an output-rich structure, but a topology in which outgoing projections are embedded in downstream targets that can continue the propagation chain.

The division-level analysis clarifies the intended scope of $\mathrm{DCA}_{\mathrm{post}}$. If cascade topology simply graded FCV continuously within every anatomical class, then controlling for brain division should have preserved a clear $\mathrm{DCA}_{\mathrm{post}}$--FCV relationship. Instead, the residual analysis indicated that the relationship was largely carried by differences among divisions, with limited additional within-division explanatory power. This pattern suggests that $\mathrm{DCA}_{\mathrm{post}}$ is best interpreted as a marker of broad dynamical regime, such as telencephalic versus non-telencephalic organization, rather than as a universal fine-grained predictor of FCV among nearby regions of the same division.

The stimulus analyses provide a functional extension of this interpretation. Regions with high spontaneous FCV also showed high OMR-associated FCV and stronger across-condition OMR FCV modulation (Fig.~\ref{fig:structure_function_synthesis}B,C). This suggests that spontaneous FCV is not merely unstable measurement or background noise, but is related to stimulus-linked modulation of a circuit's functional coupling. In this framework, telencephalic circuits may use diverse region-to-region interactions not only to maintain spontaneous dynamical flexibility, but also to support adaptive communication during changing sensory and behavioral conditions.

Cross-species analyses are consistent with broader relevance of directed cascade-like embedding, while also showing that FCV enrichment was not uniformly distributed across all functional groups. At the matched scale available in each dataset, high-FCV groups often corresponded to olfactory/chemosensory or associative/learning annotations, but these categories are species-specific and should be interpreted as functional analogies rather than homologous circuit classes. This pattern may reflect the computational role of these circuits. Olfactory/chemosensory systems must flexibly map diverse external cues onto changing internal states and behavioral responses \citep{Bargmann2006Chemosensation,Leinwand2016OlfactoryNematodes}, whereas associative/learning circuits integrate sensory, contextual, and modulatory signals to update action selection \citep{Dempsey2022ZebrafishMemorySynapses,Aso2014MushroomBodyArchitecture,Li2020MushroomBodyConnectome}. Such functions may favor circuits that can rapidly reconfigure their functional partners rather than maintain a single stable coupling pattern.

The relevant side of the directed architecture, however, was not invariant. In \textit{C.~elegans}, high-FCV olfactory/chemosensory neurons showed a $\mathrm{DCA}_{\mathrm{post}}$ tendency similar to the zebrafish output-side profile. In zebrafish, high-FCV associative/learning regions showed the clearest output-side cascade signature. This output-side pattern is consistent with anatomical work showing that zebrafish pallial zones have region-specific efferent pathways to downstream targets \citep{Yanez2022ZebrafishPalliumHodology}.In adult \textit{Drosophila}, the high-FCV mushroom body showed the opposite directional profile, with high $\mathrm{DCA}_{\mathrm{pre}}$ and low $\mathrm{DCA}_{\mathrm{post}}$. This does not necessarily contradict the zebrafish result, because the mushroom body has a different circuit architecture. It receives convergent sensory and modulatory inputs and transforms them through compartmentalized Kenyon cell, dopaminergic, and MBON pathways \citep{Aso2014MushroomBodyArchitecture,Li2020MushroomBodyConnectome}. Therefore, at the region-level scale used here, MB functional variability may be more strongly related to the directed organization of upstream inputs than to downstream output-side embedding. Thus, the cross-species comparison is consistent with directed cascade embedding as a structural correlate of FCV in high-variability sensory and associative circuits, but it does not support a single invariant DCA direction.

Direction-specific rewiring and the minimal layer model provide mechanistic plausibility for this interpretation in zebrafish (Fig.~\ref{fig:layer-linear-model}). NULL-Out disruption reduced FCV more strongly than NULL-In disruption, and increasing directional asymmetry in the layer-wise model broadened the accessible FC-state landscape. From a modeling perspective, this can be understood through the effect of asymmetric coupling on network modes. For a linearized system,
$
\dot{x}(t) = A x(t) + \eta(t)
$
symmetric coupling makes the covariance structure largely determined by the eigenvalues of $A$. When $A$ is asymmetric and non-normal, however, eigenvectors are non-orthogonal, allowing activity to be transiently amplified even when the overall system remains stable. Such transient amplification can increase the temporal variability of pairwise correlations without necessarily increasing mean coupling strength. Thus, directed cascade-like architecture may weaken stable synchronization while preserving effective signal propagation, providing a plausible route from asymmetric structural organization to higher FCV.

Several limitations qualify these conclusions. First, the zebrafish structural network should be interpreted as a synapse-informed estimate rather than a direct ultrastructural connectome. The available synapse-like endpoint data were partial and sparse relative to the full larval zebrafish brain, and synaptic directionality was inferred from soma-to-endpoint distance rather than confirmed by electron microscopy. In addition, missing local connections were supplemented using a probabilistic distance-based interpolation, with direction assigned randomly when explicit annotations were unavailable. These steps allowed us to construct a whole-brain directed SC estimate, but they may introduce inaccurate or missing edges and therefore uncertainty in regional DCA and motif estimates. DCA magnitudes may also depend on subnetwork size; threshold sensitivity and inter-regional link-composition analyses are provided in SI Appendix, Sensitivity analyses support threshold-robust inter-regional structural embedding.

Second, the zebrafish simulations provide mechanistic plausibility rather than a complete biophysical reconstruction. The whole-brain model used simplified neuronal and calcium-observation dynamics, approximate synaptic weights, and null rewiring schemes that preserved coarse anatomical region pairs while disrupting cell-level endpoints. These perturbations identify direction-specific sensitivity of the model, but they do not prove that the same manipulation would have identical effects in vivo. Cell-type identity and excitation--inhibition balance were unresolved in both the empirical SC and simulations; both likely contribute to FCV and may covary with directional embedding \citep{mueller2006phylotypic,ganz2012subpallium,Witvliet2020Connectome}. The available calcium recordings also undersample fast neural dynamics, and replication in larger cohorts will be important.

Finally, cross-species comparisons were limited by differences in functional sampling and structure--function matching. In \textit{C.~elegans}, functional recordings did not track activity from all neurons simultaneously across all animals, limiting full connectome-wide functional coverage. In \textit{Drosophila}, structure and function were matched through anatomical regions rather than one-to-one cell identity, and the structural and functional datasets were acquired from different animals and modalities. This region-level matching may obscure finer cell-level correspondences between synaptic embedding and FCV, so the fly result should be interpreted as an anatomical-scale comparison rather than a direct cell-level structure--function mapping.

In conclusion, our study shows that functional connectivity variability is associated with second-order directed cascade organization. In zebrafish, high-FCV telencephalic circuits were linked most strongly to output-side downstream embedding and output-to-output motifs. Comparative analyses in \textit{C.~elegans} and \textit{Drosophila} showed more limited directed-embedding profiles in high-FCV groups or regions, with dominant output-side or input-side expression differing across systems. Together, these findings identify second-order directed embedding as a candidate structural feature associated with FCV, strongest in the zebrafish telencephalic analysis and supported by more limited comparative patterns in two invertebrate datasets.
\matmethods{


\subsection*{Zebrafish functional activity data}

Functional activity was derived from whole-brain calcium imaging in larval zebrafish \citep{chen2018brain} (5--7 days post-fertilization; $n= 7$ animals). Imaging was acquired at approximately 2 volumes per second for $\sim$10 minutes per subject, yielding fluorescence traces from $\sim$80{,}000 neurons per animal. Signals were expressed as $\Delta F/F$ and high-pass filtered at 0.01 Hz using a zero-phase Butterworth filter. Anatomical annotation followed established zebrafish atlases as described by Mueller and Wullimann \citep{mueller2015atlas,chen2018brain}. Each neuron was assigned to an anatomical subregion label based on its registered position. Regional activity time series were obtained by averaging fluorescence signals within anatomical subdivisions; reported region-level values were obtained by averaging subdivision-level measures assigned to the same anatomical root area.

\subsection*{Synapse-informed zebrafish structural connectivity}

Directed structural connectivity was constructed by integrating neuronal soma coordinates from the calcium-imaging dataset \citep{chen2018brain} with synapse-like endpoints from a cellular-resolution anatomical dataset \citep{kunst2019cellular}. For each reconstructed neuron, terminal endpoints were identified and classified as putative presynaptic (distal) or postsynaptic (proximal) based on soma-to-endpoint distance thresholds. After aligning synaptic and neuronal coordinates, each neuron was assigned to a brain region based on its soma location. Coordinate alignment and endpoint classification are described in SI Appendix, Alignment of Synapse and Neuron Coordinates; Synaptic Classification: Pre- vs.\ Post-synaptic Endpoints; and Fig.~S1A.

Each neuron was associated with nearby synaptic sites by assigning synapses located within a Euclidean distance of $60\,\mu$m from the neuronal soma. A directed connection from neuron $i$ to neuron $j$ was defined if there existed at least one putative presynaptic site associated with neuron $i$ and at least one putative postsynaptic site associated with neuron $j$, and if both synaptic sites originated from the same ancestral neuronal morphology. In such cases, the directed adjacency element was set to $SC_{ij}= 1$. To account for potential connections not explicitly annotated, additional links were introduced probabilistically based on spatial proximity. For each unconnected pair $(i,j)$, a connection was sampled with probability decaying exponentially with Euclidean distance \citep{Kaiser2006Nonoptimal,Betzel2016Distance}, using a characteristic length scale $\lambda= 80\,\mu$m estimated from empirical distance--correlation relationships (SI Appendix, Fig.~S1B). The direction of each added connection was assigned randomly. Because directionality was inferred from anatomical endpoint geometry and random direction assignment was used only for probabilistically added links, directed structural results should be interpreted as synapse-informed estimates of projection topology rather than direct ultrastructural measurements.

\subsection*{FC and FCS}

Spontaneous functional connectivity (FC) was computed using a sliding-window Pearson correlation approach \citep{allen2014tracking,elton2015task}. For the empirical zebrafish data, 5 s windows with 2.5 s overlap were used. Within each window, pairwise correlations between regional activity time series were calculated. Functional connectivity strength (FCS) was defined as the mean absolute FC of a region with all other regions across time. Functional connectivity variability (FCV) was defined by first computing the temporal standard deviation of each pairwise FC across sliding windows and then averaging these values across all other regions of a given region. FC profile correlation-distance was computed from the temporal change in each region's whole-brain FC profile. For cross-subject comparability, region-level FC measures were $z$-scored within each animal before group-level averaging or division-level comparisons.
Region-wise FCV, FCS, FC partner reconfiguration, and TE-derived summaries are shown in SI Appendix, Fig.~S3.

\subsection*{FC partner reconfiguration}

FC partner reconfiguration quantified how strongly the pattern of a region's functional partners changed over time. For each sliding window, the FC profile of a region was defined as its set of FC values with all other regions. This profile was then compared with the region's average FC profile computed across all windows. Reconfiguration was quantified as the mean correlation distance between window-specific and average FC profiles. Higher values indicate greater changes in the identity of functional partners across time, whereas lower values indicate more stable interaction patterns. Importantly, this measure captures changes in the configuration of functional partners rather than fluctuations in overall FC strength.



\subsection*{Transfer entropy and Neighborhood transfer entropy}

Directed functional interactions were estimated using transfer entropy (TE) between community-level activity traces. Neurons were grouped into anatomical root areas, and functional communities were identified within each area using Louvain community detection. For each community, net TE drive was calculated as the average directional information transfer to other communities. Root-area values were obtained by averaging community-level measures within each anatomical region. To facilitate comparison across subjects, root-area TE values were normalized within each subject before group-level analyses. Neighborhood TE quantified the average net TE drive of functionally connected neighboring communities. Functional neighbors were identified from significant FC relationships, and neighborhood TE was computed by averaging the net TE drive of these neighboring communities. Full computational details are provided in SI Appendix, Transfer entropy and neighborhood transfer entropy; a representative TE example is shown in Fig.~S4.

\subsection*{Modularity Q and output/input ratio}

The zebrafish SC summary included modularity Q and log-transformed output/input ratio as complementary structural descriptors. Modularity Q was used to quantify local modular segregation within each anatomical region. For each subject and region, neurons assigned to that region and their within-region structural edges were used to construct a cell-level induced subnetwork. Directed edges were symmetrized, duplicate edges were combined by summed weight, and weighted community structure was estimated using a modularity objective. Subject-region modularity values were then summarized across subjects.

The output/input ratio quantified net directed projection bias. This measure was computed from inter-regional structural edges, excluding within-region edges, so that it captured macroscopic projection bias rather than local recurrent connectivity. A pseudocount of 1 was added to both outgoing and incoming strengths before taking the log ratio. Positive values indicate output-biased structural embedding, whereas negative values indicate input-biased embedding. 
Region-wise structural summaries, including DCA, modularity, log out/in ratio, and OO fraction, are shown in SI Appendix, Fig.~S2.


\subsection*{Directional Coreness Asymmetry and Output-to-Output Motif}

For each anatomical region, a local directed subnetwork $SC_k$ was extracted from the cell-level zebrafish SC matrix. To quantify directed hub embedding, we computed output and input coreness vectors for each regional directed subnetwork \citep{borgatti2000models,Borgatti2006Centrality,winding2023connectome}. The vectors $c_{\mathrm{out},k}$ and $c_{\mathrm{in},k}$ were obtained by solving
\[
(c_{\mathrm{out},k}, c_{\mathrm{in},k}) =
\arg\max_{c_{\mathrm{out}},c_{\mathrm{in}}}
c_{\mathrm{out}}^{\top} SC_k c_{\mathrm{in}}
\]
subject to
\[
c_{\mathrm{out}},c_{\mathrm{in}} \ge 0,\quad
\|c_{\mathrm{out}}\|^2 = \|c_{\mathrm{in}}\|^2 = 1 .
\]

This optimization is equivalent to rank-1 non-negative matrix factorization of the directed connectivity matrix. Numerically, continuous directional coreness was estimated using alternating rank-1 power iteration. Coreness vectors were initialized uniformly and normalized to unit Euclidean norm. At each iteration, we updated
\[
c_{\mathrm{out}}^{(t+1)}
=
\frac{SC_k c_{\mathrm{in}}^{(t)}}{\|SC_k c_{\mathrm{in}}^{(t)}\|_2},
\qquad
c_{\mathrm{in}}^{(t+1)}
=
\frac{SC_k^\top c_{\mathrm{out}}^{(t+1)}}{\|SC_k^\top c_{\mathrm{out}}^{(t+1)}\|_2},
\]
until convergence. Regions without sufficient local directed edges were assigned missing DCA values.

For each neuron $i$ in region $k$, directional coreness asymmetry (DCA) was defined as
\[
\mathrm{DCA}_i^{(k)} = c_{\mathrm{out},i}^{(k)} - c_{\mathrm{in},i}^{(k)},
\]
such that positive values indicate output-dominated embedding and negative values indicate input-dominated embedding. To characterize second-order directional topology, we defined two regional measures:
\begin{equation}
\begin{aligned}
\mathrm{DCA}_{{\mathrm{pre}},k} &=
\frac{1}{|\mathcal{P}_k|}
\sum_{i \in \mathcal{P}_k} \mathrm{DCA}_i, \\
\mathrm{DCA}_{{\mathrm{post}},k} &=
\frac{1}{|\mathcal{T}_k|}
\sum_{i \in \mathcal{T}_k} \mathrm{DCA}_i .
\end{aligned}
\end{equation}
Here, $\mathcal{P}_k$ denotes neurons projecting to region $k$, and $\mathcal{T}_k$ denotes neurons receiving projections from region $k$. High $\mathrm{DCA}_{\mathrm{post}}$ therefore indicates that a region's efferent targets are themselves output-dominated within their local circuits, whereas high $\mathrm{DCA}_{\mathrm{pre}}$ indicates that a region receives input from output-dominated upstream populations.

To validate this interpretation, we additionally computed an output-to-output (OO) inter-regional motif fraction. For each source region $k$, OO fraction was defined as the proportion of outgoing inter-regional edges for which both the source neuron and the target neuron were output-dominated within their respective local regional subnetworks:
\[
\mathrm{OO\ fraction}_k =
\frac{
\#\{i \rightarrow j:\ r(i)=k,\ r(j)\neq k,\ \mathrm{DCA}_i>0,\ \mathrm{DCA}_j>0\}
}{
\#\{i \rightarrow j:\ r(i)=k,\ r(j)\neq k\}
}.
\]
Thus, OO fraction provides a discrete motif-level estimate of broadcast-to-broadcast cascade enrichment, whereas $\mathrm{DCA}_{\mathrm{post}}$ provides a continuous embedding-based measure of the same downstream directional organization. Coreness vectors are computed independently within each regional subnetwork and normalized to unit length; DCA magnitudes therefore reflect relative directional embedding within local circuits rather than absolute broadcast strength.


%\subsection*{Region-wise SC--FC coupling analysis}

%Matched region-level SC and FC feature matrices were constructed and $z$-scored across regions. FC features included FCV, FCS, FC profile correlation-distance, $\mathrm{TE}_{\mathrm{net}}$ drive, and neighborhood $\mathrm{TE}_{\mathrm{net}}$ drive. SC features included $\mathrm{DCA}_{\mathrm{post}}$, $\mathrm{DCA}_{\mathrm{pre}}$, modularity Q, log-transformed out-/in-degree ratio, and OO fraction. Only regions with valid values in both matrices were retained.

%Pairwise associations between regional functional and structural features were assessed using Pearson correlation, with Benjamini--Hochberg FDR correction across all FC--SC feature pairs. Dominant multivariate coupling was estimated using partial least squares canonical analysis (PLS) on standardized regional FC and SC feature matrices. Specifically, FC and SC latent scores were defined as $\mathbf{t}=\mathbf{X}_{FC}\mathbf{w}_{FC}$ and $\mathbf{u}=\mathbf{X}_{SC}\mathbf{w}_{SC}$, where $\mathbf{X}_{FC}$ and $\mathbf{X}_{SC}$ denote standardized regional feature matrices. The first latent variable was used as the main SC--FC coupling axis. Regional FCV was also predicted from SC features using linear regression with SC predictors standardized within each training fold. Prediction performance was evaluated using fivefold cross-validation and reported as the observed--predicted Pearson correlation and cross-validated $R^2$.

\subsection*{Layer-wise directed linear model}

A layer-wise directed linear stochastic model was used to test whether asymmetric directed coupling can amplify FC variability. The model contained four layers with sizes $(n_1,n_2,n_3,n_4)=(2,3,4,5)$ and followed a stable linear system with diagonal decay, additive Gaussian noise, symmetric within-layer coupling, and asymmetric between-layer coupling controlled by a global asymmetry parameter $\epsilon$. In this minimal model, $\epsilon$ was used as an abstract analogue of downstream output-dominated embedding: increasing $\epsilon$ strengthened feedforward coupling toward downstream layers and weakened feedback coupling in the reverse direction. Each layer also received a slow sinusoidal drive with independently sampled phase. Dynamics were integrated with the Euler--Maruyama scheme, and FCV was estimated from sliding-window Pearson correlations of layer-averaged activity. 


For each layer $l$, neuronal activities were first averaged within the layer to obtain a layer-level signal.
For each sliding window, we computed the Pearson correlation matrix between the four layer-level signals and defined the dynamic FC state $z_l$ as the average correlation of layer $l$ with the other three layers in that window. The distribution of this state, $P_l(z)$, was estimated across windows and simulation runs, converted to an effective potential $U_l(z)=-\log P_l(z)$, and shifted so that each curve had minimum zero.
Full mathematical specification is provided in SI Appendix, Four-layer Directed Linear Network Model: Full Specification. Representative layer activity traces are shown in SI Appendix, Fig.~S8B.


\subsection*{Zebrafish whole-brain network model}

To test direction-specific effects in a synapse-informed whole-brain setting, we performed large-scale zebrafish neuronal network simulations constrained by the empirical directed structural network. Each subject-specific network contained on the order of $\sim3\times 10^{4}$ neurons and $\sim1\times 10^{6}$ directed synaptic connections. Neuronal dynamics were modeled using the Izhikevich spiking formalism \citep{izhikevich2003simple} with excitatory and inhibitory conductance-based synaptic currents and a phenomenological calcium observation model \citep{rahmati2016inferring,yaksi2006reconstruction}. Synaptic weights were initialized from empirical distance--FC relationships and refined by correlation-based optimization targeting time-averaged correlation structure, without an explicit FCV objective. FC variability was then evaluated from the fitted directed structural model and from direction-specific topology-preserving null models. Full model equations and weight optimization procedures are provided in SI Appendix, Large-scale Neuronal Simulations: Model Equations and Large-scale Network Simulation: Weight Optimization and Empirical Validation. Empirical--simulated FC validation is shown in SI Appendix, Fig.~S6.


\subsection*{Topology-preserving null models}

Direction-specific null models were constructed to test whether incoming or outgoing P/SP projections contributed more strongly to FCV. In the NULL-In condition, incoming projections to the target P/SP population were randomly rewired while preserving inter-regional connection counts. Specifically, for each incoming edge, the pre- and postsynaptic neurons were randomly reassigned within their respective anatomical regions. In the NULL-Out condition, outgoing projections from the target P/SP population were rewired analogously. These manipulations preserved coarse anatomical structure and connection density while selectively disrupting neuron-level directional embedding of incoming or outgoing P/SP projections. Null effects on DCA and representative rewired network examples are shown in SI Appendix, Fig.~S7; representative activity traces are shown in Fig.~S8A.


\subsection*{Stimulus-associated FCV}

Stimulus-associated FCV was computed from zebrafish stimulus-period calcium activity. For each subject, neuronal activity was first averaged within each anatomical region to obtain region-level stimulus traces. For each stimulus condition, time points with the same stimulus index were extracted and concatenated, and sliding-window Pearson correlation matrices were computed using 5 s windows with 2.5 s steps. Pairwise stimulus FCV was defined as the standard deviation of windowed FC values across windows, and region-level stimulus FCV was calculated as the mean pairwise FCV to all other regions. For the OMR-related analysis, forward, rightward, and leftward optomotor response conditions were retained. FCV values were $z$-scored within each subject for each stimulus condition and then averaged across subjects. Mean OMR-associated FCV was defined as the average FCV across the three OMR conditions, and across-condition OMR FCV variability was defined as the standard deviation of FCV across these three OMR conditions. Region-wise OMR FCV profiles are shown in SI Appendix, Fig.~S5.

\subsection*{\textit{C.~elegans} structural and functional analyses}

The \textit{C.~elegans} structural connectome was constructed from the adult hermaphrodite chemical synapse edge list distributed with the OpenWorm CElegansNeuroML mirror of the WormWiring/Cook 2019 connectome data \citep{Cook2019}. The Cook 2019 connectome contains 302 annotated neurons; after excluding neurons without chemical synaptic connections, 297 neurons were retained. Directed chemical synapses among annotated sensory neurons, interneurons, and motor neurons were retained, and each matrix entry represented the summed synapse weight from a presynaptic neuron to a postsynaptic neuron.

Whole-brain calcium activity data were obtained from the WormWideWeb Atanas/Kim 2023 archive \citep{AtanasKim2023}. Neuron identity labels were taken from the corresponding WormWideWeb dataset metadata and matched to structural connectome neuron names. Empty, uncertain, or alternate labels were excluded during label cleaning. Spontaneous epochs were defined as frames preceding the first annotated heat-stimulation event, and heat-epoch FCV was computed from frames following the annotated heat onset. For each recording, spontaneous calcium traces were high-pass filtered at 0.03 Hz using the native sampling rate and then $z$-scored within neuron before FC analysis. Across the 18 spontaneous recordings used for FCV analysis, 148 cleaned neuron labels were available. The final neuron-level analysis retained neurons with sufficient SC/activity matching and recording coverage.

Because the \textit{C.~elegans} connectome is defined at single-neuron resolution, directed structural embedding was computed directly on the neuron-level directed SC matrix. Output and input coreness vectors were estimated from the directed weighted SC matrix using the same directed coreness objective described above. $\mathrm{DCA}_{\mathrm{post}}$ was calculated as the outgoing synapse-weighted mean DCA of each neuron's postsynaptic targets, and $\mathrm{DCA}_{\mathrm{pre}}$ was calculated analogously from presynaptic sources. FCV was computed using sliding-window Pearson correlation with a window of 20 frames and a step of 8 frames. Within each recording, neuron-level FCV values were z-scored across neurons, and the final neuron-level zFCV value was the mean of recording-level zFCV values across recordings in which that neuron was observed.


\subsection*{\textit{Drosophila} structural and functional analyses}

The \textit{Drosophila melanogaster} structural connectome was constructed from the FlyWire783 connectome, a cell-level reconstruction of the adult female fly brain \citep{Dorkenwald2024FlyWire}. FlyWire synapse annotations were used to construct a directed cell-level SC matrix in which each edge connected one presynaptic cell to one postsynaptic cell and edge weight represented summed synapse count. Edges with fewer than five synapses were excluded. Each FlyWire cell was assigned to a side-aware Ito L/R neuropil region from synaptic localization counts. For each cell, presynaptic and postsynaptic localization counts were summed within each Ito L/R region, the region with the largest count was selected as the dominant anatomical assignment, and cells were retained only when the dominant region accounted for at least 30\% of the cell's total localization count. Here, side-aware Ito-region nodes refer to Ito anatomical neuropil labels with left and right hemispheres retained as separate nodes.

Functional activity was obtained from the Branson999 calcium imaging dataset from Turner, Mann, and Clandinin, which provides calcium recordings aligned to 999 anatomical ROIs \citep{Turner2021Branson999,Turner2021DrosophilaCentralBrainConnectivity}. Because the structural and functional datasets were acquired from different animals and modalities, they were matched through shared anatomical region labels rather than one-to-one cell identity. Each ROI was assigned to a side-aware Ito-region label\citep{Ito2014InsectBrainNomenclature}, and ROI traces sharing the same side-aware Ito region were averaged to obtain one regional activity trace per recording. These region-level traces were then linearly detrended to remove slow fluctuations, high-pass filtered at 0.03 Hz, and $z$-scored within recording before FC analysis.

Because the \textit{Drosophila} structure--function analysis was performed at anatomical-region resolution, $\mathrm{DCA}$ was computed within each Ito L/R region using the local directed cell-to-cell SC submatrix. FCV was computed from sliding-window Pearson correlations among the side-aware Ito-region traces using 15-frame windows and 5-frame steps. Node-level FCV values were $z$-scored across side-aware Ito regions within each recording and then averaged across recordings to produce the final \textit{Drosophila} FCV summary. The FCV-centered summary retained 41 side-aware Ito-region nodes with matched Branson FCV values. Structure--function analyses were restricted to the subset of nodes with complete FCV and SC features, yielding 37 matched \textit{Drosophila} nodes for the final SC--FC comparisons. 


\subsection*{Functional-group synthesis across species}

To compare the relationship between functional variability and directed structural embedding across species, we mapped species-specific neurons or brain regions onto broad comparative functional groups. These groups were not intended to define one-to-one anatomical homologies across \textit{C.~elegans}, \textit{Drosophila}, and zebrafish. Instead, they were used as coarse circuit-role categories that capture major stages of information processing: sensory input, relay or integration, sensorimotor output, and higher-order associative or state-dependent computation. This logic is consistent with comparative connectomics approaches that relate species-specific anatomical labels to broader functional or connectivity-defined axes rather than requiring exact anatomical correspondence \citep{Mars2018ConnectivityBlueprints,Xu2020CrossSpeciesAlignment,Margulies2016PrincipalGradient}. Brain regions were mapped to the following shared categories: olfactory/chemosensory, visual/optic, other sensory, sensorimotor output, integrative relay and  associative/learning.

For \textit{C.~elegans}, the classification followed the standard organization of the identified connectome into sensory neurons, interneuron or premotor-like integrative neurons, motor/command neurons, and modulatory classes \citep{Cook2019,AtanasKim2023}. Olfactory/chemosensory neurons were treated separately because chemosensation is a major sensory axis in nematode behavior and circuit organization \citep{Bargmann2006Chemosensation,Leinwand2016OlfactoryNematodes}. For \textit{Drosophila}, anatomical regions were classified using Ito-region and FlyWire anatomical labels, separating olfactory neuropils, optic neuropils, premotor/protocerebral regions, and mushroom body-related regions \citep{Ito2014InsectBrainNomenclature,Dorkenwald2024FlyWire,Turner2021Branson999}. Mushroom body regions were assigned to the associative/learning or state-related category because this circuit is a canonical center for associative learning, memory, sensory integration, and state-dependent modulation \citep{Aso2014MushroomBodyArchitecture,Li2020MushroomBodyConnectome}. For zebrafish, anatomical root areas were classified using larval brain atlas annotations and known circuit functions \citep{Wullimann1996Neuroanatomy,Randlett2015ZBrain}. Olfactory bulb and sensory regions were assigned to sensory categories, motor-associated hindbrain and spinal-projecting regions to sensorimotor output, and relay-like regions to integrative categories. Pallial and related telencephalic regions were assigned to the associative/learning or state-related category based on their roles in spatial representation, context-dependent behavior, and associative memory-related plasticity \citep{Yanez2022ZebrafishPalliumHodology,Yang2024SpatialCode,Dempsey2022ZebrafishMemorySynapses}. The detailed classification and node-level assignments are provided in SI Appendix, Table S1; zebrafish and \textit{Drosophila} anatomical region abbreviations are listed in SI Appendix, Tables S2 and S3.


\subsection*{Statistical analysis}

For zebrafish division-level FC and SC comparisons, group differences were assessed using Kruskal--Wallis tests followed by pairwise two-sided Mann--Whitney $U$ tests with Holm--Bonferroni correction. Pairwise associations between regional FC and SC features were assessed using Pearson correlation. For the zebrafish SC--FC feature matrix, Benjamini--Hochberg FDR correction was applied across all 25 FC--SC feature pairs. Dominant multivariate SC--FC coupling was estimated using partial least squares canonical analysis on standardized FC and SC feature matrices. The first latent variable was used as the main shared SC--FC axis, and its association was summarized by the correlation between FC and SC latent scores. Feature-weight stability was assessed by leave-one-region-out recomputation of the PLS model.Regional FCV prediction from SC features was evaluated using linear regression with fivefold cross-validation. Structural predictors were standardized within each training fold, and prediction performance was reported as the observed--predicted Pearson correlation and cross-validated $R^2$. Standardized regression coefficients were used to summarize the contribution of each SC feature.

For direction-specific zebrafish simulation perturbations, bootstrap distributions were generated by resampling P/SP nodes within each condition and recomputing mean FCV differences between Base, NULL-Out, and NULL-In conditions. Bootstrap confidence intervals were computed from the empirical bootstrap distributions, and bootstrap $p$ values were Holm-corrected across the reported contrasts.

For cross-species \textit{C.~elegans} and \textit{Drosophila} analyses, continuous DCA--FCV associations were assessed using Pearson correlation. For FCV group panels, global group differences were assessed using Kruskal--Wallis tests followed by pairwise two-sided Mann--Whitney $U$ tests with Holm correction. For DCA group panels, global and pairwise group differences were assessed using group-label permutation tests with Holm correction.

For functional-group synthesis, values were centered and scaled within species and metric before comparison. Functional-group enrichment was summarized as the group-versus-rest standardized mean difference. Significance was assessed by permuting functional-group labels within each species and metric while preserving group sizes, recomputing the same enrichment statistic for each permutation, and applying Holm correction across groups within each species--metric comparison.

}

\showmatmethods{} % Display the Materials and Methods section

\dataavail{All data, materials, and analysis code supporting the findings of this study are publicly available at \url{https://github.com/DongmyeongL/zebrafish_paper_project}.}

\acknow{This work was supported by the National Research Foundation of Korea (NRF) grant funded by the Korea government (MSIT) (No. 2023R1A2C2006217). R.R.\ is supported by the UK National Institute for Health and Care Research (NIHR).}



\showacknow{} % Display the acknowledgments section

%\bibsplit[3]
%Use \bibsplit to split the references from the body of the text. Value "[3]" represents the number of reference in the left column (Note: Please avoid single column figures & tables on this page.)

% Bibliography
\bibliography{reference}

\end{document}
